\documentclass{article}
\usepackage{amsmath,amssymb,amsthm}
\usepackage{algorithm}
\usepackage{algpseudocode}
\usepackage{bm}
\usepackage{enumitem}
\usepackage{hyperref}
\usepackage{geometry}
\geometry{margin=1in}
\newtheorem{lemma}{Lemma}
\newtheorem{theorem}{Theorem}
\newcommand{\EE}{\mathbb{E}}
\newcommand{\RR}{\mathbb{R}}
\newcommand{\norm}[1]{\left\lVert#1\right\rVert}
\newcommand{\inner}[1]{\left\langle#1\right\rangle}
\begin{document}

\section*{Quantized Stochastic Gradient Descent (QSGD)}

\section*{Mathematical Formulation}
Let $f:\RR^{d}\rightarrow\RR$ be the (possibly non‑convex) objective and let $\{\xi_t\}_{t\ge0}$ be an i.i.d. data stream.
At iteration $t$ we form a (stochastic) gradient
\[
g_t \;:=\; \nabla f(w_t;\xi_t)\in\RR^{d},
\]
and we communicate a \emph{quantized} version $Q(g_t)$ that satisfies  

\[
\boxed{\EE\!\left[\,Q(g_t)\mid w_t\,\right]=g_t},\qquad 
\boxed{\EE\!\left[\norm{Q(g_t)-g_t}^{2}\mid w_t\right]\le \omega\,\norm{g_t}^{2}},
\tag{1}
\]
where $\omega\ge0$ quantifies the distortion of the quantizer.

The update rule is
\[
\boxed{w_{t+1}=w_t-\eta_t\,Q(g_t)},\qquad t=0,1,\dots,T-1,
\tag{2}
\]
with a stepsize (learning rate) sequence $\{\eta_t\}_{t\ge0}$ that will be specified later.

\section*{Pseudocode}
\begin{algorithm}[H]
\caption{Quantized Stochastic Gradient Descent (QSGD)}
\begin{algorithmic}[1]
\State \textbf{Input:} initial point $w_0\in\RR^{d}$, stepsize schedule $\{\eta_t\}$, quantizer $Q(\cdot)$
\For{$t=0,1,\dots,T-1$}
    \State Sample a mini‑batch (or a datapoint) $\xi_t$
    \State Compute stochastic gradient $g_t=\nabla f(w_t;\xi_t)$
    \State Quantize: $\tilde g_t = Q(g_t)$   \Comment{Unbiased by (1)}
    \State Update: $w_{t+1}=w_t-\eta_t\,\tilde g_t$   \Comment{Equation (2)}
\EndFor
\State \textbf{Output:} $\{w_t\}_{t=0}^{T}$ (or the averaged iterate $\bar w_T:=\tfrac1T\sum_{t=0}^{T-1}w_t$)
\end{algorithmic}
\end{algorithm}

\section*{Dry Run Example}
Consider the scalar quadratic $f(w)=(w-3)^{2}$, smooth with $L=2$.
We use the identity quantizer $Q(g)=g$ (hence $\omega=0$), stepsize $\eta=0.1$, and start at $w_0=0$.

\begin{center}
\begin{tabular}{c|c|c|c}
$t$ & $g_t=\nabla f(w_t)$ & $w_{t+1}=w_t-\eta g_t$ & $f(w_{t+1})$\\\hline
0 & $2(w_0-3) = -6$ & $0-0.1(-6)=0.6$ & $(0.6-3)^2=5.76$\\
1 & $2(0.6-3) = -4.8$ & $0.6-0.1(-4.8)=1.08$ & $(1.08-3)^2=3.6864$\\
2 & $2(1.08-3) = -3.84$ & $1.08-0.1(-3.84)=1.464$ & $(1.464-3)^2=2.3660$
\end{tabular}
\end{center}
The objective value decreases at each step, illustrating the basic behaviour of QSGD.

\newpage
\section*{Assumptions}
We work under the following standard conditions for non‑convex stochastic optimisation with unbiased quantisation.

\begin{enumerate}[label=(A\arabic*)]
\item \textbf{Smoothness.} $f$ is $L$‑smooth: for all $x,y\in\RR^{d}$
\[
f(y)\le f(x)+\inner{\nabla f(x)}{y-x}+\frac{L}{2}\norm{y-x}^{2}.
\tag{3}
\]

\item \textbf{Unbiased Stochastic Gradient with Bounded Variance.}
\[
\EE\!\left[g_t\mid w_t\right]=\nabla f(w_t),\qquad 
\EE\!\left[\norm{g_t-\nabla f(w_t)}^{2}\mid w_t\right]\le\sigma^{2}.
\tag{4}
\]

\item \textbf{Unbiased Quantiser with Bounded Relative Error.}
Equation \eqref{1} holds for a constant $\omega\ge0$ (e.g. $\omega= \frac{d}{s}-1$ for $s$‑bit stochastic rounding).

\item \textbf{Stepsize.} For the main theorem we choose a constant stepsize
\[
\eta_t\equiv\eta\le\frac{1}{L(1+\omega)}.
\tag{5}
\]
A diminishing schedule $\eta_t=\frac{\eta_0}{\sqrt{t+1}}$ is also admissible (see discussion).
\end{enumerate}

\section*{Main Convergence Theorem}
\begin{theorem}[Non‑convex convergence of QSGD]\label{thm:main}
Assume (A1)–(A4) and let the constant stepsize $\eta$ satisfy \eqref{5}.  
Define the average squared gradient norm
\[
\Phi_T\;:=\;\frac{1}{T}\sum_{t=0}^{T-1}\EE\!\left[\norm{\nabla f(w_t)}^{2}\right].
\]
Then
\[
\boxed{\Phi_T\;\le\;
\frac{2\bigl(f(w_0)-f^{*}\bigr)}{\eta T}
\;+\;
L\eta\,(1+\omega)\,\sigma^{2}},
\tag{6}
\]
where $f^{*}\triangleq\inf_{w}f(w)$.
Consequently, choosing $\eta=\Theta\!\bigl(1/\sqrt{T}\bigr)$ yields
\[
\Phi_T = O\!\Bigl(\frac{1}{\sqrt{T}}\Bigr).
\]
\end{theorem}

\section*{Theoretical Properties}
\begin{itemize}
\item \textbf{Stability.} The condition $\eta\le 1/(L(1+\omega))$ guarantees that the term
$1-\frac{L\eta}{2}(1+\omega)$ in the descent inequality stays non‑negative, preventing divergence caused by quantisation noise.
\item \textbf{Hyper‑parameter Sensitivity.} The bound \eqref{6} shows a linear dependence on $\omega$; finer quantisation (smaller $\omega$) directly improves the asymptotic rate.
\item \textbf{Comparison to Baselines.}
    \begin{itemize}
    \item With $\omega=0$ (exact gradients) \eqref{6} reduces to the classic SGD bound for non‑convex smooth objectives.
    \item For $\omega>0$ the additional term $L\eta\omega\sigma^{2}$ quantifies the penalty incurred by compression.
    \end{itemize}
\end{itemize}

\section*{Discussion}
The theorem matches the optimal $O(1/\sqrt{T})$ rate for non‑convex stochastic optimisation up to the multiplicative factor $1+\omega$ caused by quantisation.  
If a diminishing stepsize $\eta_t=\eta_0/\sqrt{t+1}$ is used, the same $O(1/\sqrt{T})$ rate is obtained without requiring the explicit upper bound \eqref{5} (the proof is identical after replacing $\eta$ by $\eta_t$ and summing a harmonic series).

\newpage
\section*{Preliminaries (Definitions and Lemmas)}
\begin{lemma}[Smoothness inequality]\label{lem:smooth}
If $f$ is $L$‑smooth, then for any $x,y\in\RR^{d}$
\[
f(y)\le f(x)+\inner{\nabla f(x)}{y-x}+\frac{L}{2}\norm{y-x}^{2}.
\]
\end{lemma}
\begin{lemma}[Unbiasedness of the composite estimator]\label{lem:unbiased}
For any $t$,
\[
\EE\!\left[Q(g_t)\mid w_t\right]=\EE\!\left[\EE\!\left[Q(g_t)\mid g_t,w_t\right]\mid w_t\right]
      =\EE\!\left[g_t\mid w_t\right]=\nabla f(w_t).
\]
\end{lemma}
\begin{lemma}[Second‑moment bound for the quantised estimator]\label{lem:second}
Using (1) and (4),
\[
\EE\!\left[\norm{Q(g_t)}^{2}\mid w_t\right]
\le (1+\omega)\Bigl(\norm{\nabla f(w_t)}^{2}+\sigma^{2}\Bigr).
\tag{7}
\]
\end{lemma}
\begin{proof}
\begin{align*}
\EE\!\left[\norm{Q(g_t)}^{2}\mid w_t\right]
 &=\EE\!\left[\norm{g_t+(Q(g_t)-g_t)}^{2}\mid w_t\right]\\
 &=\EE\!\left[\norm{g_t}^{2}\mid w_t\right]
   +2\EE\!\left[\inner{g_t}{Q(g_t)-g_t}\mid w_t\right]
   +\EE\!\left[\norm{Q(g_t)-g_t}^{2}\mid w_t\right]\\
 &\overset{(a)}{=}\EE\!\left[\norm{g_t}^{2}\mid w_t\right]
   +\EE\!\left[\norm{Q(g_t)-g_t}^{2}\mid w_t\right] \qquad\text{(unbiasedness)}\\
 &\le\EE\!\left[\norm{g_t}^{2}\mid w_t\right]+\omega\,\EE\!\left[\norm{g_t}^{2}\mid w_t\right]\\
 &= (1+\omega)\EE\!\left[\norm{g_t}^{2}\mid w_t\right]\\
 &\overset{(b)}{\le}(1+\omega)\bigl(\norm{\nabla f(w_t)}^{2}+\sigma^{2}\bigr).
\end{align*}
Step (a) follows from Lemma~\ref{lem:unbiased}; step (b) uses the variance bound (4). ∎
\end{proof}

\section*{Main Proof (Step‑by‑Step Derivation)}
We prove Theorem~\ref{thm:main}. All expectations are taken with respect to the joint randomness of $\{\xi_t\}$ and the quantiser.

\begin{enumerate}[label=[Step \arabic*], leftmargin=*, itemsep=4pt]
\item \textbf{Apply smoothness.}  
Using Lemma~\ref{lem:smooth} with $x=w_t$, $y=w_{t+1}=w_t-\eta Q(g_t)$,
\[
f(w_{t+1})\le f(w_t)-\eta\inner{\nabla f(w_t)}{Q(g_t)}+\frac{L\eta^{2}}{2}\norm{Q(g_t)}^{2}.
\tag{8}
\]

\item \textbf{Take conditional expectation w.r.t. $w_t$.}  
Denote $\mathcal{F}_t=\sigma(w_0,\dots,w_t)$. Taking $\EE[\cdot\mid\mathcal{F}_t]$ on both sides of \eqref{8} yields
\[
\EE\!\left[f(w_{t+1})\mid\mathcal{F}_t\right]
\le f(w_t)-\eta\,\EE\!\left[\inner{\nabla f(w_t)}{Q(g_t)}\mid\mathcal{F}_t\right]
      +\frac{L\eta^{2}}{2}\EE\!\left[\norm{Q(g_t)}^{2}\mid\mathcal{F}_t\right].
\tag{9}
\]

\item \textbf{Handle the inner product term.}  
By Lemma~\ref{lem:unbiased},
\[
\EE\!\left[\inner{\nabla f(w_t)}{Q(g_t)}\mid\mathcal{F}_t\right]
 =\inner{\nabla f(w_t)}{\EE[Q(g_t)\mid\mathcal{F}_t]}
 =\inner{\nabla f(w_t)}{\nabla f(w_t)}
 =\norm{\nabla f(w_t)}^{2}.
\tag{10}
\]

\item \textbf{Bound the second‑moment term.}  
Lemma~\ref{lem:second} gives
\[
\EE\!\left[\norm{Q(g_t)}^{2}\mid\mathcal{F}_t\right]
\le (1+\omega)\bigl(\norm{\nabla f(w_t)}^{2}+\sigma^{2}\bigr).
\tag{11}
\]

\item \textbf{Insert (10) and (11) into (9).}  
\[
\EE\!\left[f(w_{t+1})\mid\mathcal{F}_t\right]
\le f(w_t)-\eta\norm{\nabla f(w_t)}^{2}
   +\frac{L\eta^{2}}{2}(1+\omega)\bigl(\norm{\nabla f(w_t)}^{2}+\sigma^{2}\bigr).
\tag{12}
\]

\item \textbf{Collect the gradient‑norm terms.}  
Rearranging \eqref{12},
\[
\EE\!\left[f(w_{t+1})\mid\mathcal{F}_t\right]
\le f(w_t)-\underbrace{\eta\Bigl(1-\frac{L\eta}{2}(1+\omega)\Bigr)}_{\displaystyle \alpha}
   \norm{\nabla f(w_t)}^{2}
   +\frac{L\eta^{2}}{2}(1+\omega)\sigma^{2}.
\tag{13}
\]

\item \textbf{Choose stepsize to keep $\alpha>0$.}  
Assumption (A4) guarantees $\eta\le 1/(L(1+\omega))$, hence
\[
1-\frac{L\eta}{2}(1+\omega)\;\ge\;\frac12,
\]
so that $\alpha\ge \frac{\eta}{2}>0$.

\item \textbf{Take full expectation and sum over $t=0,\dots,T-1$.}  
Denote $F_t:=\EE[f(w_t)]$. From \eqref{13} and the tower property,
\[
F_{t+1}\le F_t-\alpha\,\EE\!\left[\norm{\nabla f(w_t)}^{2}\right]
          +\frac{L\eta^{2}}{2}(1+\omega)\sigma^{2}.
\tag{14}
\]
Summing \eqref{14} for $t=0$ to $T-1$ telescopes the $F_t$ terms:
\[
F_T\le F_0-\alpha\sum_{t=0}^{T-1}\EE\!\left[\norm{\nabla f(w_t)}^{2}\right]
      +\frac{L\eta^{2}}{2}(1+\omega)\sigma^{2}T.
\tag{15}
\]

\item \textbf{Isolate the average gradient norm.}  
Since $F_T\ge f^{*}$,
\[
\alpha\sum_{t=0}^{T-1}\EE\!\left[\norm{\nabla f(w_t)}^{2}\right]
\le F_0-f^{*}
    +\frac{L\eta^{2}}{2}(1+\omega)\sigma^{2}T.
\tag{16}
\]
Dividing by $\alpha T$ and using $\alpha\ge\eta/2$,
\[
\frac{1}{T}\sum_{t=0}^{T-1}\EE\!\left[\norm{\nabla f(w_t)}^{2}\right]
\le \frac{2(F_0-f^{*})}{\eta T}
   +L\eta(1+\omega)\sigma^{2}.
\tag{17}
\]

\item \textbf{Conclude.}  
Equation \eqref{17} is exactly the bound \eqref{6} claimed in Theorem~\ref{thm:main}. ∎
\end{enumerate}

\section*{Rate Derivation (With Constants)}
From \eqref{6} we have
\[
\Phi_T\le \frac{2\Delta}{\eta T}+L\eta(1+\omega)\sigma^{2},
\qquad \Delta:=f(w_0)-f^{*}.
\]
Choosing $\eta=\sqrt{\frac{2\Delta}{L(1+\omega)\sigma^{2}T}}$ (which satisfies \eqref{5} for sufficiently large $T$) balances the two terms and yields
\[
\Phi_T\le 2\sqrt{\frac{2L(1+\omega)\Delta\sigma^{2}}{T}}
      =O\!\bigl(T^{-1/2}\bigr).
\]

\section*{Special Cases}
\begin{description}
\item[Exact gradients ($\omega=0$)] The bound reduces to the classic SGD result.
\item[Deterministic setting ($\sigma=0$)] The second term vanishes and we obtain a $O(1/T)$ decay of the gradient norm.
\item[Vanishing stepsize $\eta_t=\eta_0/\sqrt{t+1}$] Repeating the above proof with $\eta_t$ yields
\[
\frac{1}{T}\sum_{t=0}^{T-1}\EE\!\left[\norm{\nabla f(w_t)}^{2}\right]
\le \frac{C}{\sqrt{T}},
\]
where $C$ depends on $L,\omega,\sigma,\eta_0$ and $\Delta$.
\end{description}

\end{document}